\section{Reduced and thick boundary contact}

We isolate the cancellation calculation because it supplies the numerical
invariant used in the degreewise theorem.  The first proof is global and
eliminative; the second is local and computes the contact order without a
resultant.

\begin{proposition}[all-parameter cancellation trace]
\label{prop:thick-trace}
For every $m,r\ge1$, let $\delta_{m,r}$ be the reduced prime discriminant
equation fixed in Section~4.  Up to multiplication by a nonzero scalar,
\[
 \delta_{m,r}(0,Q,R)
 =Q^{mr(m+1)}\bigl((r+1)RQ^m-C\bigr).                    \tag{7.1}
\]
Consequently
\[
 \frac{\C[Q,R]}{(\delta_{m,r}(0,Q,R))}
 \simeq
 \frac{\C[Q,R]}{(Q^{mr(m+1)})}
 \times\C[Q,Q^{-1}],                                    \tag{7.2}
\]
and the nilradical has exact nilpotency index $mr(m+1)$.
\end{proposition}

\begin{proof}[First proof: monic resultant and subresultant]
On the normalized critical divisor put $Y=Q-PT$.  Then
\[
 T=Y^{-m},\quad P=(Q-Y)Y^m.                              \tag{7.3}
\]
After substituting $t=uY^{-m}$ in the critical-value equation and clearing
the unit $Y^{-m(r+1)}$, define
\[
 F(Y)=P-(Q-Y)Y^m,                                        \tag{7.4}
\]
\[
 G(Y)=C I_{m,r}(Y,Q)-RY^{m(r+1)},                        \tag{7.5}
\]
where
\[
I_{m,r}(Y,Q)=
 \int_0^1\left[Y^m-u\{Q+(Y-Q)u\}^m\right]^r\,du.        \tag{7.6}
\]
Equivalently, the integral-free coefficient formula is
\[
 I_{m,r}=
 \sum_{j=0}^r\sum_{k=0}^{mj}
 \frac{(-1)^j}{j+k+1}\binom rj\binom{mj}k
 Y^{m(r-j)}Q^{mj-k}(Y-Q)^k.
\]
Thus (7.6) is a polynomial with rational coefficients.  The monic
polynomial $F$ is primitive over $\mathbb Z[C,P,Q,R]$.  Multiplication of
$G$ by $(mr+r+1)!$ clears every integration denominator; taking its primitive
part gives an integral polynomial with the same zero scheme over $\C$.  This
works uniformly for every symbolic pair of positive integers $(m,r)$.  It is
the precise uniform integral meaning here: $m$ and $r$ occur as exponents,
so there is no literal polynomial ring in indeterminates $m,r$ containing all
members at once.

The polynomial $F$ is monic in $Y$.  Its incidence with $G$ is the integral
critical-value incidence of Section~4 and maps birationally to the reduced
discriminant.  At its generic point, $F$ and $G$ have exactly one common
simple root: the critical points are simple and their critical values are
distinct.  Thus the zeroth subresultant specializes to zero and the first
subresultant does not, certifying a gcd of degree one.  More concretely,
after adjoining the distinct roots of
$F$, the resultant is the product of the $G$-values; generically exactly one
factor vanishes, and it is transverse because $G$ is linear in $R$ with
nonzero coefficient at that root.  It follows that the primitive resultant
\[
 D_{m,r}(P,Q,R)=\operatorname{Res}_Y(F,G)                 \tag{7.7}
\]
is a nonzero scalar multiple of the reduced prime generator
$\delta_{m,r}$, not a power of it.

Now specialize $P=0$.  Since $F(Y)=Y^m(Y-Q)$ is monic,
\[
 D_{m,r}(0,Q,R)=G(0)^mG(Q).                              \tag{7.8}
\]
The two values of (7.6) are uniform beta integrals:
\[
 I_{m,r}(0,Q)
 =(-1)^rQ^{mr}\frac{r!(mr)!}{(mr+r+1)!},                 \tag{7.9}
\]
\[
 I_{m,r}(Q,Q)=\frac{Q^{mr}}{r+1}.                        \tag{7.10}
\]
Hence $G(0)$ is a nonzero scalar times $CQ^{mr}$, while
\[
 G(Q)=Q^{mr}\left(\frac{C}{r+1}-RQ^m\right).            \tag{7.11}
\]
Substitution in (7.8) gives (7.1).  This derivation takes the resultant only
after $F,G$ have been made polynomial and primitive; no component can be
lost through denominator clearing.

The two factors in (7.1) are comaximal because the second is congruent to
$-C$ modulo $Q$.  The Chinese remainder theorem gives (7.2).  In the first
factor the nilradical is generated by $Q$ and has exact index $mr(m+1)$; the
second factor is reduced.
\end{proof}

\begin{proof}[Second proof: completed local contact]
This proof computes the exponent in (7.1) without using a resultant.  Let
$\eta$ be the generic point of the affine-line component $P=Q=0$ of the
intersection, and put $K=\C(R)$.  The completed target local ring transverse
to that line is
\[
 \widehat{\mathcal O}_{\A^3,\eta}=K[[P,Q]].              \tag{7.12}
\]
Use the completed regular parameter plane $A=K[[Y,Q]]$ for the polynomial
critical-value incidence; its slice over the generic value $R$ is $G=0$.
The pullback of the plane $P=0$ is the Cartier divisor
\[
 P=(Q-Y)Y^m,
 \qquad
 \operatorname{div}(P)=m[Y=0]+[Y=Q].                    \tag{7.13}
\]
Neither $Y$ nor $Q-Y$ divides $G$, by (7.9)--(7.10).  Additivity of
intersection multiplicity for the divisor (7.13) therefore gives
\[
 \operatorname{length}_K A/(P,G)
 =m\,\operatorname{ord}_QG(0,Q)
  +\operatorname{ord}_QG(Q,Q).                           \tag{7.14}
\]
Equations (7.9) and (7.11), now used only as completed-local restrictions,
give
\[
 \operatorname{ord}_QG(0,Q)=mr,
 \quad
 \operatorname{ord}_QG(Q,Q)=mr,                         \tag{7.15}
\]
because $C/(r+1)-RQ^m$ is a unit in $K[[Q]]$.  Consequently
\[
 \operatorname{length}_K A/(P,G)
 =m(mr)+mr=mr(m+1).                                      \tag{7.16}
\]

The critical-value normalization is birational with residue degree one.
The local projection formula therefore identifies (7.16) with the order of
$\delta_{m,r}(0,Q,R)$ in $K[[Q]]$.  Thus the completed intersection ring at
$\eta$ is
\[
 K[[Q]]/(Q^{mr(m+1)}),                                   \tag{7.17}
\]
which directly gives contact order and nilpotency index $mr(m+1)$.  The two
summands in (7.14) also explain the exponent geometrically: $m^2r$ comes
from the multiplicity-$m$ branch $Y=0$, and $mr$ from the transverse branch
$Y=Q$.
\end{proof}

\begin{lemma}[contact calculation]\label{lem:contact}
For the split weighted map, $Z_\Delta\cap Z_0$ is reduced and $\mu=1$.
For a cancellation map of type $(m,r)$ with $N\ge4$, it is nonreduced and
\[
 e_\Delta=r+1,\qquad
 \mu=mr(m+1)=m(N-1)>1.                                   \tag{7.18}
\]
\end{lemma}

\begin{proof}
For the weighted seed, Lemma~\ref{lem:weighted-saturation} gives
\[
 C=0,\qquad \delta_N|_{C=0}=\mu_N(B^2-8A)=0,
\]
with $\mu_N\ne0$.  Its coordinate ring is $\C[B]$, hence is reduced.

For cancellation, $e_\Delta=r+1$ was computed in Lemma~\ref{lem:boundary},
and Proposition~\ref{prop:thick-trace} gives the scheme intersection and its
exact nilpotency index.  In the range $N\ge4$, $P=0$ is a canonical boundary
vertex by Lemma~\ref{lem:boundary}, and the index exceeds one.
\end{proof}
